

\section{Sobre el tipo de homotopía y el pegado de $\lambda$-células \color{blue}(título provisional)}

En esta sección sentaremos las bases de la Teoría de Morse. Hasta ahora hemos aportado el aparataje técnico sobre el que construir la teoría, aquí veremos como las funciones de Morse nos pueden servir para reconstruir una variedad mediante el pegado de $\lambda$-células, donde $\lambda$ representa el índice de un punto crítico no degenerado.

Antes de probar los probar los resultados principales de esta sección conviene mencionar qué entendemos por \textit{variedad riemanniana} y dar una definición de campo vectorial.

\begin{definition}
Sea $M$ una variedad diferenciable.  
Llamamos \textit{fibrado tangente} de $M$ a la unión disjunta
de los espacios tangentes en todos los puntos de $M$, es decir,
\[
TM \;=\; \bigsqcup_{p \in M} T_pM.
\]
\end{definition}

\begin{definition}
Un \textit{campo vectorial suave} sobre una variedad diferenciable $M$ es una aplicación suave
\[
X \colon M \to TM,
\]
que denotamos $p \mapsto X_p$, y que verifica que, para todo $p \in M$, se tiene $X_p \in T_pM$.
\end{definition}

\begin{definition}
Una \textit{variedad riemanniana} es un par $(M,\langle\cdot,\cdot\rangle)$ donde $M$ es una variedad diferenciable y para cada punto $p \in M$, $\langle\cdot,\cdot\rangle_p$ es un producto escalar (luego, bilineal, simétrico y definido positivo)
  \[
    \langle\cdot,\cdot\rangle_p \colon T_pM \times T_pM \to \mathbb{R},
  \]
que verifica que para cualesquiera campos vectoriales suaves $X$ e $Y$, la función
    \[
      p \mapsto \langle X(p), Y(p) \rangle_p
    \]
    es suave sobre $M$.
\end{definition}


\bigskip
\begin{definition}
    Sea $M$ una variedad diferenciable. Llamamos \textit{grupo uniparamétrico de difeomorfismos de M} a una aplicación $\varphi:\mathbb{R}\times M \to M$ que verifica:
    \begin{enumerate}[label=\roman*)]
        \item $\forall t\in \mathbb{R},~~ \varphi_t:M\to M$ definida como $\varphi_t(q)=\varphi(t,q)$, es un difeomorfismo de $M$ sobre si misma.

        \item $\forall t, s\in\mathbb{R}$ se tiene que $\varphi_{t+s}=\varphi_t\circ \varphi_s$.
    \end{enumerate}
\end{definition}

\bigskip

Dado un grupo uniparamétrico de difeormorfismos $\varphi$ definimos el siguiente campo vectorial sobre $M$

\begin{equation*}
    X_q(f)=\lim_{h\to 0}\frac{f(\varphi_h(q))-f(q)}{h}
\end{equation*}
donde $f$ es una función suave de $M$ en $\mathbb{R}$, y decimos que $X$ genera $\varphi$.

Normalmente uno se refiere a la curva $\varphi_\bullet(q): \mathbb{R}\to M,~~ t\mapsto \varphi_t(q)$ como el \textit{flujo de X} cuando se conoce el campo vectorial $X$ que genera el grupo uniparamétrico de difeomorfismos. En cualquier caso, mantendremos un vocabulario lo más fiel posible al texto \cite{milnor1963}.

También hablaremos del \textit{vector velocidad} de una curva diferenciable $c:\mathbb{R}\to M$ que definimos de manera similar a X
\begin{equation*}
    \frac{dc}{dt}(f)= \lim_{h\to 0}\frac{f(c(t+h))-f(c(t))}{h}\in T_{c(t)}M 
\end{equation*} 

\bigskip
\begin{lemma}
\label{unidiff}
    Todo campo vectorial sobre una variedad diferenciable $M$ que se anula fuera de un compacto $K\subset M$ genera un único grupo de difeomorfismos de M.
\end{lemma}
{\color{blue} Esta demostración se me hace liosa y no se si está del todo bien o si se entiende lo suficientemente bien.}
\begin{proof}
    Sean $X$ un campo vectorial y $\varphi$ un grupo uniparamétrico de difeomorfismos de $M$ generado por $X$. Entonces se tiene que fijado un punto cualquiera $q\in M$ la curva $\varphi_\bullet(q): \mathbb{R}\to M,~~ t\mapsto \varphi_t(q)$ verifica la ecuación diferencial
    \begin{equation*}
        \frac{d\varphi_t}{dt}(q)=X_{\varphi_t(q)}\tag{$\ast$}
    \end{equation*}
    con condición inicial $\varphi_0(q)=q$.\\

    En efecto,
    \begin{align*}
        \frac{d\varphi_t(q)}{dt}(f)=\lim_{h\to 0} \frac{f(\varphi_{t+h}(q))- f(\varphi_t(q))}{h} = \lim_{h\to 0} \frac{f(\varphi_{h}(\varphi_t(q)))- f(\varphi_t(q))}{h} = \lim_{h\to 0} \frac{f(\varphi_{h}(p))- f(p)}{h} = X_p(f)
    \end{align*}

    con $p=\varphi_t(q)$.\\

    Ahora, por el Teorema de Picard-Lindelöf sabemos que existe localmente una solución única, que depende de manera diferenciable de la condición inicial. Es decir, $\forall p\in M$ existen un entorno $\mathcal{U}$ de $p$ y $\varepsilon > 0$ tales que $(\ast)$ tiene una única solución para todo $q\in \mathcal{U}$ y $|t|<\varepsilon$.\\

    Como $K\subset M$ es compacto, puede ser cubierto por una cantidad finita de dichos entornos. Sea $\varepsilon_0$ el menor de los valores definidos arriba y establezcamos $\varphi_t(q)=q,~~\forall q\notin K$ (luego $X_q=0$ fuera de $K$). Entonces $(\ast)$ tiene una única solución $\varphi_t(q)$ para $|t|<\varepsilon_0$ y $\forall q\in M$ suave respecto de ambas variables.\\

    Veamos que para todo $t$ adecuado cada $\varphi_t$ es un difeomorfismo de $M$ en si misma. Para ello basta comprobar que se cumple que $\varphi_{t+s}=\varphi_t\circ \varphi_s$ suponiendo que $|t|,|s|,|t+s|<\varepsilon_0$.\\

    Sean $q\in M$ fijo y  $p=\varphi_s(q)$, entonces $\varphi_t(p)=\varphi_t(\varphi_s(q))$. Definimos la curva $\gamma(u)=\varphi_{u+s}(q)$, así, $\gamma(0)=p$ y 
    
    \begin{equation*}
        \frac{d\gamma}{du}=\frac{d\varphi_{u+s}}{du}(q)=X_{\varphi_{u+s}(q)}=X_{\gamma (u)}
    \end{equation*} 
    
    Por lo que $\gamma$ es también una solución de ($\ast$) con condición inicial $\gamma(0)=p$.\\

    Sin embargo, tambien $\varphi_u(p)$ es solución con condición inicial $\varphi_0(p)=p$. Luego por unicidad de las soluciones
    \begin{equation*}
        \varphi_{u+s}(q)=\gamma(u)=\varphi_u(p)=\varphi_u(\varphi_s(q))=\varphi_u\circ \varphi_s(q)
    \end{equation*}
    Siempre que $|u|, |u+s|<\varepsilon_0$.

    Por último, necesitamos definir $\varphi_t$ para todo $t\in \mathbb{R}$.\\

    Sea $t\in \mathbb{R}$. Sean $k\in \mathbb{Z}$ y $r\in \mathbb{R}$ con $|r|\le \frac{\varepsilon_0}{2}$ los únicos valores tales que $t=k\frac{\varepsilon_0}{2}+r$. Entonces definimos

    \begin{equation*}
        \varphi_t=\underbrace{\varphi_{\frac{\varepsilon_0}{2}}\circ \dots \circ \varphi_{\frac{\varepsilon_0}{2}}}_{k~veces}\circ \varphi_r 
    \end{equation*}

    Es claro que $\varphi:\mathbb{R}\times M\to M$ es un difeomorfismo único para todo $t\in \mathbb{R}$.

\end{proof}

\begin{definition}
    Sean $(M,\langle\cdot,\cdot\rangle)$ una variedad riemanniana y $f:M\to\mathbb{R}$ suave. Se define el \textit{gradiente de f} como el campo vectorial $\nabla f$ que verifica 
    \begin{equation*}
        \langle X,\nabla f\rangle = X(f)
    \end{equation*}
    para cualquier otro campo vectorial $X$.
\end{definition}
\bigskip
% aqui deberia escribir algo para introducir lo siguiente.

\begin{definition}
    Sean $M$ una variedad diferenciable y $f:M\to\mathbb{R}$ suave definimos
    \begin{equation*}
        M^a := f^{-1}(-\infty,a] =\{p\in M: f(p)\le a\}
    \end{equation*}
\end{definition}

\begin{theorem}
\label{retractodef}
    Sea $f:M\to\mathbb{R}$ suave sobre una variedad riemanniana $(M,\langle\cdot,\cdot\rangle)$. Sean $a<b$ reales tales que $f^{-1}[a,b]$ es compacto y no contiene ningún punto crítico de $f$. Entonces $M^a$ es difeomorfo a $M^b$.\\
    Además, $M^a$ es retracto de deformación de $M^b$, luego tienen en el mismo tipo de homotopía.
\end{theorem}

\begin{proof}
    Sea $c:\mathbb{R}\to M$ una curva con vector velocidad $\frac{dc}{dt}$, entonces
    \begin{equation*}
        \langle \frac{dc}{dt}, \nabla f\rangle = \frac{dc}{dt}(f)=\frac{d(f\circ c)}{dt}
    \end{equation*}

    Definimos la aplicación suave $\rho : M\to \mathbb{R}$ como

    \begin{equation*}
        \rho(p):=
        \begin{cases}
            \frac{1}{\langle \nabla_pf, \nabla_pf\rangle}~~,~p\in f^{-1}[a,b]\\
            0 ~~~~~~~~~~~~~~,p\notin K
        \end{cases}
    \end{equation*}

    donde $K$ es un entorno compacto de $f^{-1}[a,b]$\\

    Ahora, sea el campo vectorial dado por $X_q=\rho(q)\nabla_pf$. Se tiene que $X_p$ se anula fuera de $K$ luego por el Lema \ref{unidiff} genera un único grupo uniparamétrico de difeomorfismos $\varphi_t$ de $M$.\\

    Fijemos $q\in M$ y consideremos $f\circ\varphi_{\bullet}(q),~t\mapsto f(\varphi_t(q))$. Si $\varphi_t(q)$ está en $f^{-1}[a,b]$ se tiene
    \begin{align*}
        \frac{df(\varphi_t(q))}{dt} = \left\langle \frac{d\varphi_t(q)}{dt},\nabla f\right\rangle = \langle X, \nabla f\rangle = \langle \rho~ \nabla f, \nabla f\rangle =\\ =\left\langle \frac{1}{\langle \nabla f,\nabla f\rangle}\nabla f,\nabla f\right\rangle=
        \frac{\langle \nabla f,\nabla f\rangle}{\langle \nabla f,\nabla f\rangle}= 1
    \end{align*}

    por lo que $f(\varphi_t(q))$ es lineal con derivada $1$ siempre que $f(\varphi_t(q))$ esté en $[a,b]$. Es decir, $f(\varphi_t(q))=f(q)+t$ {\color{blue}Esto tengo la duda de que sea exactamente así, pero funciona bien.}\\

    Veamos ahora que $\varphi_{b-a}$ es un difeomorfismo de $M^a$ en $M^b$.\\

    Sea $q\in M$,
    \begin{equation*}
        \varphi_{b-a}(q)\in M^b \iff(\varphi_{b-a}(q))=f(q)+b-a\le b \iff f(q)-a\le 0 \iff q\in M^a
    \end{equation*}
     y en concreto manda la frontera de $M^a$ a la frontera de $M^b$
     \begin{equation*}
         q\in \partial M^a \iff f(q)=a \implies f(\varphi_{b-a}(q))=b \iff b\in \partial M^b.
     \end{equation*}

     Para concluir la demostración falta comprobar que efectivamente $M^a$ es retractor de deformación de $M^b$. Para ello vamos a construir una homotopía relativa a $M^a$ de la identidad de $M^b$ a una retracción de $M^b$ a $M^a$.
     
    Sea $H(t,q):[0,1]\times M^b\to M^b$ definida como
    \begin{equation*}
        H(t,q)=
        \begin{cases}
            q~~~~~~~~~~~~~~~~~,f(q)\le a\iff q\in M^a\\
            \varphi_{t(a-f(q))}(q)~~,a\le f(q)\le b
        \end{cases}
    \end{equation*}

    Entonces, $H(0,\cdot) = id_{M^b}$ ya que $\varphi_0(q)=q,~\forall q\in M^b$ y $H(1,\cdot)$ es una retracción de $M^b$ en $M^a$, comprobémoslo. Sea $q\in f^{-1}[a,b]$, se tiene
    \begin{equation*}
        H(1,q)=f(\varphi_{(a-f(q))}(q))= f(q)+(a-f(q))= a \iff \varphi_{(a-f(q))}(q)\in f^{-1}(a)
    \end{equation*}

    luego $H(1,H(1,\cdot))= id_{M^a}$.    

\end{proof}
\bigskip

\begin{theorem}
     Sea $f:M\to\mathbb{R}$ suave sobre una variedad riemanniana $(M,\langle\cdot,\cdot\rangle)$ y $p\in M$ un punto crítico no degenerado de $f$ con índice $\lambda$.\\
     Supongamos que $f(p)=c$ y que  $\exists\varepsilon>0$ tal que $f^{-1}[c-\varepsilon,c+\varepsilon]$ es compacto y no contiene ningún punto crítico de $f$ a parte de $p$. Entonces, para todo $0<\varepsilon_0<\varepsilon$ suficientemente pequeño, $M^{c+\varepsilon}$ tiene el tipo de homotopía de $M^{c-\varepsilon}$ con una $\lambda$-célula pegada.
\end{theorem}

\begin{proof}
    Vamos a definir una función $F$ que sea menor que $f$ en un entorno $\mathcal{U}$ de $p$ con el fin de dejar a $p$ fuera de $F^{-1}[c-\varepsilon, c+\varepsilon]$ y poder aplicar a $F$ el Teorema \ref{retractodef}. Así, habremos demostrado que $F^{-1}(-\infty,c-\varepsilon]=M^{c-\varepsilon}_F$ es retracto de deformación de $M^{c+\varepsilon}_f$.
    
    Para finalizar la demostración probaremos que $M^{c-\varepsilon}_f$ con una $\lambda$-célula pegada es retracto de deformacion de $M^{c-\varepsilon}_F$ y por la transitividad de esta relación habremos terminado.\\
%Preliminares------------------

    Sean una carta $(\mathcal{U}, \phi)$ de $M$ que contiene a $p$ y $(u^1,\dots , u^n)$ las coordenadas inducidas tal que $u^1(p)=\dots=u^{n}(p)=0$. Con el fin de simplificar las expresiones que siguen definimos las funciones $\xi:\mathcal{U}\to [0,\infty)$ y $\eta:\mathcal{U}\to [0,\infty)$ como
    \begin{gather*}
        \xi(q)=(u^1(q))^2+\dots +(u^{\lambda}(q))^2\\
        \eta(q)=(u^{\lambda +1}(q))^2+\dots +(u^{n}(q))^2
    \end{gather*}
    Aplicando el Lema de Morse (\ref{lemamorse}) a $f$ obtenemos

    \begin{equation*}
        f(q) = c - \xi(q) + \eta(q)
    \end{equation*}
    $\forall q\in \mathcal{U}$, donde $c=f(p).$\\

    Sea $0<\varepsilon_0\le\varepsilon$ suficientemente pequeño para que:
    \begin{enumerate}[label=\roman*)]
        \item $f^{-1}[c-\varepsilon_0,c+\varepsilon_0]$ sea compacto y no contenga ningún punto crítico de $f$ distinto de $p$.
        \item $\phi^{-1}(\mathcal{U})\subset \mathbb{R}^n$ contenga a la bola cerrada
        \begin{equation*}
            \phi\left(\overline{B(0, \sqrt{2\varepsilon_0})}\right)= \{q\in \mathcal{U}: \xi(q)+\eta(q)\le 2\varepsilon_0\}
        \end{equation*}
    \end{enumerate}

%Construcción de F----------------
    Construyamos $F:M\to \mathbb{R}$. Sea $\mu:\mathbb{R}\to\mathbb{R}$ una aplicación suave que verifica lo siguiente
    
    \begin{equation*}
        \begin{cases}
        \mu(0)>\varepsilon_0\\
        \mu(r)=0, r\ge 2\varepsilon_0\\
        -1<\mu'(r)\le 0, \forall r\in \mathbb{R}
        \end{cases}
    \end{equation*}
    entonces definimos $F$ como
    
    \begin{equation*}
        F(q)=
        \begin{cases}
        f(q)~~~~~~~~~~~~~~~~~~~~~~~~,~~q\notin \mathcal{U}\\
        f(q)- \mu(\xi(q)+2\eta(q)),~~q\in \mathcal{U}.
        \end{cases}
    \end{equation*}
    Es claro que $F$ esta bien definida ya que $\mu$ se anula en $\partial\mathcal{U}$. Además, es suave ya que en un entorno de $\partial\mathcal{U}$ se tiene que $\mu'$ es identicamente $0$, luego tanto aproximándonos desde dentro de $\mathcal{U}$ como desde fuera $d_qF=d_qf$.\\
    
%Primera retraccion-------------------------------

    Ahora que tenemos $F$ debemos comprobar que nos encontramos en las hipótesis del Teorema \ref{retractodef}. Veamos primero que $M_F^{c+\varepsilon_0}=M_f^{c+\varepsilon_0}$, después que $F$ y $f$ poseen los mismos puntos críticos y que $F^{-1}[c-\varepsilon_0,c+\varepsilon_0]$ no contiene ningún punto crítico.\\

    En efecto, $\mu(\xi(q)+2\eta(q))=0, ~\forall q\in M$ tal que $\xi(q)+2\eta(q)\ge 2\varepsilon_0$ luego fuera del elipsoide $E=\{q\in M:\xi(q)+2\eta(q)\le 2\varepsilon_0\}$ se tiene que $F=f$.\\
    
    Sabemos que $E\in M_f^{c+\varepsilon_0}$ y por
    \begin{equation*}
        F(q)<f(q)=c-\xi(q)+\eta(q)\le c+\frac{1}{2}\xi(q)+\eta\le c + \varepsilon_0, ~\forall q\in E
    \end{equation*}
    se tiene que $E\in M_F^{c+\varepsilon_0}$. Visto que es la única región en la que difieren concluimos que $M_F^{c+\varepsilon_0}=M_f^{c+\varepsilon_0}$.\\

    Ahora veamos que coinciden sus puntos críticos. Observamos que en $\mathcal{U}$
    \begin{align*}
        \frac{\partial F}{\partial \xi}&=-1-\mu'(\xi + 2\eta)<0\text{ y}\\
        \frac{\partial F}{\partial \eta}&=1-2\mu'(\xi+2\eta)\ge 1
    \end{align*}
    Los puntos críticos de $F$ serán los que anulen su derivadas parciales
    \begin{align*}
        \frac{\partial F}{\partial x_i}=\frac{\partial F}{\partial \xi}\frac{\partial \xi}{\partial x_i} + \frac{\partial F}{\partial \eta}\frac{\partial \eta}{\partial x_i}
    \end{align*}
    pero el único punto que verifica simultáneamente $\frac{\partial \xi}{\partial x_i}=0$ y $\frac{\partial \eta}{\partial x_i}=0$ es $p$ ya que si no $f$ tendría más puntos críticos en $\mathcal{U}$. Luego como fuera de $\mathcal{U}$ coinciden se tiene que  $F$ y $f$ tienen los mismos puntos críticos.\\

    Veamos que $F^{-1}[c-\varepsilon_0,c+\varepsilon_0]=M^{c+\varepsilon_0}_F\setminus M^{c-\varepsilon_0}_F$ es compacto y no contiene ningún punto crítico. Es claro que $M^{c-\varepsilon_0}_f\subset M^{c-\varepsilon_0}_F$ ya que $F$ y $f$ coinciden salvo en el entorno de $p$ en el que $F\le f$ {\color{blue}explicar esto mejor?} así
    \begin{equation}
        M^{c+\varepsilon_0}_F=M^{c+\varepsilon_0}_f\implies M^{c+\varepsilon_0}_F\setminus M^{c-\varepsilon_0}_F \subset M^{c+\varepsilon_0}_f\setminus M^{c-\varepsilon_0}_f=f^{-1}[c-\varepsilon_0,c+\varepsilon_0]
    \end{equation}
    por lo que $F^{-1}[c-\varepsilon_0,c+\varepsilon_0]$ es compacto. Por otro lado,
    \begin{equation*}
        F(p)=c-\mu(0)<c-\varepsilon_0
    \end{equation*}
    luego $p\notin F^{-1}[c-\varepsilon_0,c+\varepsilon_0]$.\\

    De esta manera, nos encontramos en las hipótesis del Teorema \ref{retractodef} para $F$, lo que demuestra que $M^{c-\varepsilon_0}_F$ es retracto de deformación de $M^{c+\varepsilon_0}_F=M^{c+\varepsilon_0}_f$.\\

%Segunda retracción------------------

    Sea $e^{\lambda}=\{q\in \mathcal{U}: \xi(q)\le \varepsilon_0,~\eta(q)=0 \}\subset \mathcal{U}$ una $\lambda$-célula. En primer lugar, veamos que
    \begin{equation*}
        \partial e^{\lambda}=\{q\in \mathcal{U}: \xi(q)= \varepsilon_0,~~\eta(q)=0 \}=M^{c-\varepsilon_0}_f\cap e^{\lambda}
    \end{equation*}

    \begin{quote}
    $\left( \subseteq \right) ~~\forall q\in \partial e^{\lambda},~f(q)=c-\varepsilon_0 \implies q\in M^{c-\varepsilon_0}_f$\\
    
    $\left( \supseteq \right) ~~\forall q\in M^{c-\varepsilon_0}_f\cap e^{\lambda}$, se tiene $\xi(q)\le\varepsilon_0,~\eta(q)=0 \text{ y } f(q)\le -\varepsilon_0$, entonces como
    \begin{equation*}
        f(q)=c-\xi(q)\le c-\varepsilon_0\iff \xi(q)\ge \varepsilon_0
    \end{equation*}
    \quad \quad necesariamente $\xi(q)=\varepsilon_0$, luego $q\in \partial e^{\lambda}$. 
    \end{quote}
    \bigskip

    De cara al resto de la demostración resultará conveniente introducir la noción de \textit{asa}. Entenderemos $M_F^{c-\varepsilon_0}$ como $M_f^{c-\varepsilon_0}\cup H$ donde $H=\overline{M_F^{c-\varepsilon_0}\setminus M_f^{c-\varepsilon_0}}$ es el asa del que hablamos, denotado así en virtud de su nombre en inglés \textit{handle}.\\
    
    Observamos que $e^{\lambda}$ está contenido en $H$. En efecto, $F$ alcanza su máximo en $\mathcal{U}$ cuando $\mu(\xi+2\eta)$ alcanza su mínimo y este lo alcanza cuando $\xi$ y $\eta$ se anulan, es decir, únicamente en $p$, así
    \begin{equation*}
        F(q)\le F(p)<c-\varepsilon_0 ~\text{ y }~f(q)=c-\xi(q)\ge c-\varepsilon_0,~\forall q\in e^{\lambda}
    \end{equation*}
    luego $q\in e^{\lambda}\implies q\in H$.\\

    
    Ahora queremos demostrar que $M^{c-\varepsilon_0}_f\cup e^{\lambda}$ es retracto de deformación de $M^{c-\varepsilon_0}_F = M_f^{c-\varepsilon_0}\cup H$. Para ello vamos a construir una homotopía relativa a $M^{c-\varepsilon_0}_f\cup e^{\lambda}$ de identidad a una retracción, como en la demostración del Teorema \ref{retractodef}.\\

    Sea $R:[0,1]\times M_f^{c-\varepsilon_0}\cup H\to M_f^{c-\varepsilon_0}\cup H$, definimos $D$ como la identidad fuera de $\mathcal{U}$. En $\mathcal{U}$ debemos distinguir tres regiones.\\

    \textbf{Región 1:}\quad $R_1=\{q\in\mathcal{U}:\xi(q)\le \varepsilon_0\}$.\\
    
    Definimos $D(t,q)=D(t,(u_1(q),\dots,u_n(q))=(u_1(q),\dots,u_{\lambda}(q),tu_{\lambda+1}(q),\dots,tu_n(q))$.\\
    
    Es claro que es relativa a $M^{c-\varepsilon_0}_f\cup e^{\lambda}$ y   $D(1,\cdot)$ es la identidad. Además,
    \begin{equation*}
        \forall q\in R_1, ~D(0, q)\in e^{\lambda} ~\text{ y }~ D(0,D(0,q))=D(0,q).
    \end{equation*} 

    Asimismo, como $\frac{\partial F}{\partial \eta}>0$ (es decir, $F$ crece en la dirección de $\eta$) se tiene que $F(D(t,q))\le F(q)$ para todo $t\in [0,1]$ y todo $q\in R_1$. Luego $D$ lleva $M_f^{c-\varepsilon_0}\cup H$ en si mismo.


    \textbf{Región 2:}\quad $R_2=\{q\in \mathcal{U}:~ \varepsilon_0\le \xi(q)\le \eta(q)+\varepsilon_0\}$.\\

    Definimos $D(t,q)=D(t,(u_1(q),\dots,u_n(q))=(u_1(q),\dots,u_{\lambda}(q),s_t(q) u_{\lambda+1}(q),\dots,s_t(q) u_n(q))$, donde 
    \begin{equation*}
        s_t(q)=t + (1-t)\left(\frac{\xi(q)-\varepsilon_0}{\eta}\right)^{1/2}
    \end{equation*}

    

    
    

    

    
\end{proof}

{\color{blue} Al definir $\xi,\eta$ estamos dando un nuevo sistema de coordenadas, y por tanto, una nueva base del tangente. Podríamos llamarlas coordenadas de Morse. En el caso 1 de la segunda retraccion estamos contrayendo los puntos en la direccion del eje $\eta$ (vertical en los diagramas), esto es, en la direccion del campo vectorial $\frac{\partial}{\partial\eta}$...}
    